%% capitoli/cap4_procedimento.tex
%% In testa al file, prima di scrivere, incollare l'estratto di
%% output/F4-indice.md relativo a questo capitolo (tesi sostenuta,
%% fonti, lunghezza stimata): serve a controllare gli scostamenti.
%% Convenzione delle etichette: sez:cap4:<nome>, tab:cap4:<nome>.

\chapter{Il nuovo procedimento consultivo}
\label{cap:cap4}

\section{Legittimazione soggettiva e competenza dell'organo}
\label{sez:cap4:legittimazione}


\section{L'oggetto: questione interpretativa su fattispecie concreta, soglia di valore, esclusioni}
\label{sez:cap4:oggetto}


\section{Il termine e la disciplina del silenzio}
\label{sez:cap4:termine}


\section{Istruttoria, motivazione, pubblicità}
\label{sez:cap4:istruttoria}


\section{La nomofilachia delle Sezioni riunite}
\label{sez:cap4:nomofilachia}


